%% ma: L11-B14-0020
%% lop: 11
%% chuong: 4
%% bai: 14
%% dang: 11.14.4
%% loai: TN4
%% muc: VD
%% dap_an: "A"
%% nguon: "tai-lieu-hoc-tap-mon-toan-lop-11-theo-chuong-trinh-giao-duc-pho-thong-2018.pdf (Lớp Toán Thầy Duy, Phạm Lê Duy, THPT Chu Văn An), Chương 4 Bài 5, Phần 1, Câu 15 (tr.529)"
%% kien_thuc: "Lớp 11 Bài 14 (phép chiếu song song)"
%% trang_thai: chinh-thuc
%% ghi_chu: "Lời giải dùng tọa độ afin để kiểm; cách giải thuần túy (giao tuyến của (AMN) với (SBC), (SCD)) cần giáo viên bổ sung. Đã duyệt 03/11/2026."
\begin{ex}
Cho hình chóp $S.ABCD$ có đáy $ABCD$ là hình thang với $AB\parallel CD$, $AB=3CD$. Gọi $M,N$ theo thứ tự là trung điểm $SB,SC$ và $K$ là giao điểm của đường thẳng $SD$ với mặt phẳng $(AMN)$. Tính tỷ số $\dfrac{SK}{SD}$.
\choice{\True $\dfrac{SK}{SD}=\dfrac35$.}{$\dfrac{SK}{SD}=\dfrac12$.}{$\dfrac{SK}{SD}=\dfrac47$.}{$\dfrac{SK}{SD}=\dfrac23$.}
\loigiai{\sloppy Chọn tọa độ $A(0;0;0)$, $B(3;0;0)$, $D(0;1;0)$, $C(1;1;0)$ (thỏa $AB\parallel CD$, $AB=3CD$) và đỉnh $S$ bất kỳ. Mặt phẳng $(AMN)$ với $M,N$ trung điểm của $SB,SC$ cắt $SD$ tại $K$ thỏa $\overrightarrow{SK}=k\overrightarrow{SD}$. Giải điều kiện $K\in(AMN)$ được $k=\dfrac35$ (kiểm bằng sympy với $S$ tùy ý; tỉ số không phụ thuộc $S$). Vậy $\dfrac{SK}{SD}=\dfrac35$.}
\end{ex}
